\section{Integration Approach}
\label{sec:integration_approach}

This chapter describes the integration strategy employed during the development of the Best Bike Paths system, detailing the bottom-up approach in the backend, parallel frontend development, and the final synchronization phase.

\subsection{Overall Strategy}
\label{subsec:overall_strategy}

The integration process followed a \textbf{hybrid approach}:

\begin{itemize}
    \item \textbf{Backend:} Bottom-Up integration starting from the data layer
    \item \textbf{Frontend:} Component-based development with incremental feature integration
    \item \textbf{Synchronization:} Regular integration points where frontend and backend were connected and tested together
\end{itemize}

This strategy provided several advantages:
\begin{itemize}
    \item \textbf{Parallel Development:} Frontend and backend teams could work independently
    \item \textbf{Early Validation:} Core backend components were validated before API integration
    \item \textbf{Reduced Risk:} Issues were identified within individual layers before full system integration
    \item \textbf{Iterative Testing:} Tests were written during integration and refined as components matured
\end{itemize}

\section{Backend Integration (Bottom-Up)}
\label{sec:backend_integration}

The backend followed a strict bottom-up integration approach, building from foundational layers to higher-level components.

\subsection{Phase 1: Data Layer Foundation}
\label{subsec:phase1_data_layer}

\textbf{Objective:} Establish a stable data persistence layer before implementing business logic.

\textbf{Components Integrated:}
\begin{itemize}
    \item PostgreSQL database with PostGIS extension
    \item Database schema creation via SQL migrations
    \item Connection pool configuration
    \item Repository pattern implementation (User, Trip, Path, Rating repositories)
    \item Redis client for caching and session storage
\end{itemize}

\textbf{Integration Activities:}
\begin{enumerate}
    \item Set up Docker Compose infrastructure (PostgreSQL, Redis, RabbitMQ)
    \item Run database migrations to create schema
    \item Implement repository modules with CRUD operations
    \item Test repository methods against live database instances
    \item Verify referential integrity and constraint enforcement
\end{enumerate}

\textbf{Testing During Integration:}
\begin{itemize}
    \item Repository tests with mocked database pool
    \item SQL query validation
\end{itemize}

\subsection{Phase 2: Service Layer Integration}
\label{subsec:phase2_service_layer}

\textbf{Objective:} Implement business logic that orchestrates repository operations.

\textbf{Components Integrated:}
\begin{itemize}
    \item Authentication Service (registration, login, JWT generation)
    \item User Service (profile management)
    \item Trip Service (manual trip recording, statistics calculation)
    \item Path Service (CRUD operations, state transitions)
    \item Weather Service (external API integration with fallback)
    \item Email Service (SMTP configuration with mock mode)
\end{itemize}

\textbf{Integration Activities:}
\begin{enumerate}
    \item Connect services to repositories
    \item Implement service-to-service communication (e.g., Trip Service calls Weather Service)
    \item Configure RabbitMQ message queues for asynchronous tasks
    \item Implement retry logic and error handling
    \item Test services with mocked repositories
\end{enumerate}

\textbf{Testing During Integration:}
\begin{itemize}
    \item Unit tests for service logic (authService, pathService, tripService, weatherService)
    \item Integration tests for service-repository interactions
    \item Mock external API responses
    \item Validate business rules (e.g., password hashing, token expiration)
\end{itemize}

\subsection{Phase 3: API Layer Integration}
\label{subsec:phase3_api_layer}

\textbf{Objective:} Expose services through RESTful API endpoints.

\textbf{Components Integrated:}
\begin{itemize}
    \item Express.js application setup
    \item Controller layer (Auth, Trip, Path controllers)
    \item Route definitions (/api/v1/auth, /api/v1/trips, /api/v1/paths)
    \item Authentication middleware (JWT verification)
    \item Request validation middleware
    \item Error handling middleware
\end{itemize}

\textbf{Integration Activities:}
\begin{enumerate}
    \item Create controllers that call services
    \item Define API routes and HTTP methods
    \item Implement authentication middleware chain
    \item Add request body validation
    \item Configure CORS for frontend communication
    \item Standardize error response format
\end{enumerate}

\textbf{Testing During Integration:}
\begin{itemize}
    \item API route tests using Supertest
    \item Authentication flow tests (401, 403 responses)
    \item HTTP method validation (GET, POST, PUT, DELETE)
    \item Request/response payload validation
\end{itemize}

\subsection{Phase 4: Infrastructure Integration}
\label{subsec:phase4_infrastructure}

\textbf{Objective:} Integrate message queuing and caching for production-ready features.

\textbf{Components Integrated:}
\begin{itemize}
    \item RabbitMQ consumer for path statistics updates
    \item Email queue for asynchronous notifications
\end{itemize}

\textbf{Integration Activities:}
\begin{enumerate}
    \item Configure RabbitMQ exchanges and queues
    \item Implement consumer services
    \item Test message publishing and consumption
\end{enumerate}

\textbf{Testing During Integration:}
\begin{itemize}
    \item Infrastructure tests (server health check, Redis connection)
    \item Message queue tests (publishing, consumer processing)
\end{itemize}

\section{Frontend Development (Parallel)}
\label{sec:frontend_development}

While the backend followed bottom-up integration, the frontend was developed in parallel using a component-based approach.

\subsection{Phase 1: Core Components \& Navigation}
\label{subsec:frontend_phase1}

\textbf{Objective:} Establish UI foundation and navigation structure.

\textbf{Components Developed:}
\begin{itemize}
    \item Reusable UI components (Button, Card, Input, Loading)
    \item Navigation stack and tab structure
    \item Authentication context for state management
    \item Splash screen and home screen
\end{itemize}

\textbf{Testing:}
\begin{itemize}
    \item Component unit tests (Button.test.js, Card.test.js, Input.test.js)
    \item Navigation flow tests
    \item Context provider tests
\end{itemize}

\subsection{Phase 2: Authentication Screens}
\label{subsec:frontend_phase2}

\textbf{Objective:} Implement user registration and login flows.

\textbf{Components Developed:}
\begin{itemize}
    \item LoginScreen with form validation
    \item RegisterScreen with password strength checks
    \item ProfileScreen for user info display
    \item API client configuration with axios
    \item Session storage using AsyncStorage
\end{itemize}

\textbf{Testing:}
\begin{itemize}
    \item Screen rendering tests (LoginScreen.test.js, RegisterScreen.test.js)
    \item E2E authentication flow (authFlow.test.js)
    \item API mock tests for login/register endpoints
\end{itemize}

\subsection{Phase 3: Path Management Screens}
\label{subsec:frontend_phase3}

\textbf{Objective:} Implement path browsing, viewing, and creation.

\textbf{Components Developed:}
\begin{itemize}
    \item PathListScreen for browsing published paths
    \item PathDetailScreen for viewing path information
    \item PathCreateManualScreen for manual path creation
    \item PublishPathScreen for publishing paths
    \item PathCard and PathList components
    \item Map integration with Mapbox
\end{itemize}

\textbf{Testing:}
\begin{itemize}
    \item Path component tests (PathCard.test.js, PathList.test.js)
    \item E2E path flow (pathFlow.test.js)
    \item Map service tests (MapService.test.js)
\end{itemize}

\subsection{Phase 4: Offline Support \& Sync}
\label{subsec:frontend_phase4}

\textbf{Objective:} Add offline capabilities and data synchronization.

\textbf{Components Developed:}
\begin{itemize}
    \item Network status detection
    \item Offline banner components
    \item Sync queue for pending operations
    \item Storage service wrappers
    \item Retry logic for failed requests
\end{itemize}

\textbf{Testing:}
\begin{itemize}
    \item Storage tests (storageService.test.js, syncQueue.test.js)
    \item Offline sync tests (syncOffline.test.js)
    \item Network status tests
\end{itemize}

\section{Frontend-Backend Synchronization}
\label{sec:synchronization}

After developing frontend and backend components in parallel, several integration and synchronization sprints ensured seamless communication.

\subsection{API Contract Definition}
\label{subsec:api_contract}

\textbf{Activities:}
\begin{itemize}
    \item Documented API endpoints, request/response schemas
    \item Aligned frontend API client with backend route definitions
    \item Standardized error response format
    \item Agreed on authentication header format (Bearer token)
\end{itemize}

\subsection{Integration Testing Sessions}
\label{subsec:integration_sessions}

\textbf{Sprint 1: Authentication Integration}
\begin{itemize}
    \item Connected LoginScreen to POST /api/v1/auth/login
    \item Connected RegisterScreen to POST /api/v1/auth/register
    \item Verified token storage and retrieval
    \item Tested session expiration and logout
\end{itemize}

\textbf{Sprint 2: Path Management Integration}
\begin{itemize}
    \item Connected PathListScreen to GET /api/v1/paths
    \item Connected PathDetailScreen to GET /api/v1/paths/:id
    \item Connected PathCreateManualScreen to POST /api/v1/paths
    \item Verified map rendering with backend geospatial data
\end{itemize}

\textbf{Sprint 3: Trip Recording Integration}
\begin{itemize}
    \item Connected trip recording to POST /api/v1/trips
    \item Integrated weather data display
    \item Verified trip statistics calculation
    \item Tested trip history retrieval
\end{itemize}

\subsection{End-to-End Validation}
\label{subsec:e2e_validation}

\textbf{Complete User Flows Tested:}
\begin{enumerate}
    \item \textbf{UC1 - Registration:} New user creates account → Receives verification email → Logs in
    \item \textbf{UC2 - Login:} User logs in → Token stored → Authenticated requests succeed
    \item \textbf{UC3 - Manual Path Creation:} User creates path → Saves as draft → Publishes → Appears in search
    \item \textbf{UC6 - Path Search:} User searches paths → Views results → Selects path → Views details
    \item \textbf{UC8 - Trip History:} User records trip → Views trip in history → Sees statistics
\end{enumerate}

\section{Testing Methodology}
\label{sec:testing_methodology}

Tests were implemented \textbf{during integration} and \textbf{refined after integration} as the system matured.

\subsection{Test-Driven Integration}
\label{subsec:test_driven_integration}

The integration process followed these practices:

\begin{itemize}
    \item \textbf{Write Repository Tests First:} Before implementing service logic, repository tests ensured data layer correctness
    \item \textbf{Mock Lower Layers:} Service tests mocked repositories; controller tests mocked services
    \item \textbf{Integration Tests for Boundaries:} Tested service-repository and controller-service boundaries
    \item \textbf{E2E Tests Last:} Full user flows tested after all components integrated
\end{itemize}

\subsection{Test Evolution}
\label{subsec:test_evolution}

As the system evolved, tests were refined:

\begin{itemize}
    \item \textbf{Initial Tests:} Simple happy-path scenarios
    \item \textbf{Mid-Integration:} Added error handling tests (404, 401, 500 responses)
    \item \textbf{Post-Integration:} Edge cases, boundary conditions, concurrent operations
    \item \textbf{Final Refinement:} Performance tests, load tests, security tests (SQL injection, XSS)
\end{itemize}

\subsection{Continuous Integration}
\label{subsec:ci_pipeline}

A CI pipeline automated testing throughout integration:

\begin{itemize}
    \item \textbf{On Commit:} Linting and unit tests run
    \item \textbf{On Pull Request:} Full integration test suite runs with Dockerized infrastructure
    \item \textbf{On Merge:} Deployment to staging environment for manual QA
\end{itemize}

\section{Integration Challenges \& Solutions}
\label{sec:integration_challenges}

\subsection{Challenge 1: Mapbox API integration in expo}
\label{subsec:challenge_mapbox_expo}

\textbf{Issue:} Expo does not support Mapbox API by itself, so to avoid any issues, we decided to keep it just for the calculations. For the map we later used the react-native-maps API.

\subsection{Challenge 2: Git branching and merging}
\label{subsec:challenge_git_merging}

\textbf{Issue:} During the integration phase, we had to simultaneously work on the same files, which led to some merge conflicts that significantly slowed down the integration process.

\section{Integration Tools \& Environment}
\label{sec:integration_tools}

\subsection{Testing Tools}
\label{subsec:testing_tools}

\textbf{Backend:}
\begin{itemize}
    \item \textbf{Jest:} Test framework for unit and integration tests
    \item \textbf{Supertest:} HTTP assertion library for API testing
    \item \textbf{Docker:} Containerized infrastructure for consistent test environments
\end{itemize}

\textbf{Frontend:}
\begin{itemize}
    \item \textbf{Jest:} Unit and integration testing
    \item \textbf{React Native Testing Library:} Component rendering and interaction simulation
    \item \textbf{Expo Go:} Device emulation for manual testing
\end{itemize}

\subsection{Environment Configuration}
\label{subsec:environment_config}

The integration environment uses Docker Compose:

\begin{verbatim}
docker-compose.yml        # PostgreSQL, Redis, RabbitMQ
docker-compose.dev.yml    # Backend service with hot-reload
\end{verbatim}

Environment variables control:
\begin{itemize}
    \item Database connection strings
    \item API keys (mocked in test environment)
    \item Feature flags
    \item Log levels
\end{itemize}

\subsection{Test Data Management}
\label{subsec:test_data_management}

Test data is managed to ensure reproducibility:
\begin{itemize}
    \item \textbf{Fixtures:} JSON files define static test data
    \item \textbf{Factories:} Functions generate randomized test entities
    \item \textbf{Cleanup:} Database transactions rolled back or tables truncated between tests
\end{itemize}
